% ============================================================================
% 07_casos_uso_reales.tex — Cuerpo del entregable "Casos de Uso Reales:
% diseño de interfaz gráfica". Formato fully dressed (Larman) a nivel real
% (concreto), apoyado en las pantallas y callouts del Entregable 04.
% Sin preámbulo: se incluye tanto desde 07_Casos_Uso_Reales.tex (standalone)
% como, eventualmente, desde un documento maestro.
% ============================================================================
\section{Casos de Uso Reales: diseño de interfaz gráfica}

\subsection{Introducción}

Esta sección redacta, en formato completamente vestido (\textit{fully
dressed}) según Craig Larman, la versión \textbf{real} (concreta) de los
casos de uso principales de LSA Salud: describe la interacción efectiva con la
interfaz gráfica real del sistema, en lugar de la intención abstracta del
actor.
Cada caso de uso real toma como base su contraparte \textbf{esencial},
documentada en el Entregable~03 (Modelos de Casos de Uso), y la especializa
citando pantallas (P-01 a P-05) y elementos numerados (\emph{callouts})
del Entregable~04 (Mockups); por ejemplo, ``el sistema muestra la palabra
reconocida en el panel (6) de la Pantalla P-02''.

Se desarrollan aquí CU-01 (Comunicarse mediante Lengua de Señas Argentina),
CU-04 (Responder al Paciente) con su extensión CU-05 (Dictar Respuesta por
Voz), CU-06 (Capturar Muestra de Seña para Entrenamiento), CU-09 (Otorgar
Consentimiento Explícito de Cámara), CU-10 (Reportar Traducción Incorrecta
o Ambigua) y CU-12 (Finalizar Atención y Preparar la Estación para el
Próximo Paciente). El escenario principal de éxito de cada caso de uso real
se presenta en el formato de dos columnas propuesto por Larman, con las
acciones del actor a la izquierda y las respuestas del sistema a la
derecha. Los IDs de actores, casos de uso y requisitos citados son los
mismos que en el resto de los entregables.

% ============================================================================
\subsection{CU-01 (real): Comunicarse mediante Lengua de Señas Argentina}

Especializa el caso de uso esencial \textbf{CU-01} (ver Entregable~03,
sección 2.3) sobre la Pantalla \textbf{P-02 --- Estación del Paciente}
(Figura~\ref{fig:cur-p02-cu01}).

\diagrama{04_mockup_p02_paciente}{\textwidth}{Pantalla P-02, referenciada por CU-01 (real) y CU-04 (real)}{fig:cur-p02-cu01}

\begin{longtable}{@{}p{4.1cm} >{\raggedright\arraybackslash}p{11.2cm}@{}}
\toprule
\textbf{Nombre} & Comunicarse mediante Lengua de Señas Argentina (real) \\
\textbf{ID} & CU-01 (real; esencial en Entregable 03) \\
\textbf{Actor principal} & Paciente \\
\textbf{Personal involucrado e intereses} &
Paciente: quiere expresar una necesidad médica sin depender de un
intérprete humano. \newline
Personal de Recepción o Enfermería: necesita recibir el mensaje del
paciente en español, de forma clara y oportuna, en su propia pantalla (P-03). \\
\textbf{Precondiciones} &
El paciente abrió la Pantalla P-02 (\texttt{/paciente}); el navegador
autorizó el acceso a la cámara; el indicador (3) muestra ``EN VIVO''. \\
\textbf{Garantías de éxito} &
El mensaje traducido aparece como una burbuja en el panel de conversación
(7) de P-02 y en el panel (3) de P-03 (CU-04, real), y se reprodujo el
audio si la síntesis estuvo disponible. \\
\bottomrule
\caption{CU-01 (real) — Comunicarse mediante Lengua de Señas Argentina}
\end{longtable}

\setlength{\extrarowheight}{5pt}
\begin{longtable}{@{}>{\raggedright\arraybackslash}p{0.47\textwidth} >{\raggedright\arraybackslash}p{0.47\textwidth}@{}}
\toprule
\textbf{Acción del actor} & \textbf{Respuesta del sistema} \\
\midrule
\endfirsthead
\toprule
\textbf{Acción del actor} & \textbf{Respuesta del sistema} \\
\midrule
\endhead
1. El paciente se ubica frente a la cámara, visible en el área de video
(4) de la Pantalla P-02. &
2. El sistema dibuja en vivo los puntos de referencia (\textit{landmarks})
sobre la imagen del paciente en el área de video (4). \\
3. El paciente hace una seña. &
4. El sistema actualiza en tiempo real las tres barras del panel
``Top-3 Predicciones'' (6) con la seña reconocida y su porcentaje de
confianza. \\
&
5. Cuando la confianza de la seña ganadora supera el 70\,\% y no es la
seña de reposo (``idle''), el sistema agrega la palabra a la ventana de
oración en construcción (5), superpuesta sobre el video. \\
6. El paciente encadena varias señas. &
7. El sistema va sumando palabras a la ventana (5) (por ejemplo, ``yo
cabeza dolor''). \\
8. El paciente baja las manos y deja de moverse. &
9. El sistema detecta la quietud, cierra la secuencia de señas y
reemplaza el contenido de la ventana (5) por el indicador ``Generando
respuesta de IA...''. \\
&
10. El sistema traduce la secuencia a una oración natural (CU-02, real) y
agrega automáticamente una nueva burbuja al panel de conversación (7) de
P-02, con el texto traducido (por ejemplo, ``Me duele la cabeza.''); la
ventana (5) desaparece, sin que medie ningún botón de confirmación entre
los pasos 9 y 10. \\
&
11. El sistema sintetiza el audio de la oración (CU-03, real) y lo
reproduce automáticamente, sin acción del paciente. \\
&
12. La misma burbuja aparece, al mismo tiempo, en el panel (3) de la
Pantalla P-03 del personal (ver CU-04, real). \\
\bottomrule
\caption{Curso normal de los eventos --- CU-01 (real)}
\end{longtable}
\setlength{\extrarowheight}{0pt}

\begin{longtable}{@{}p{4.1cm} >{\raggedright\arraybackslash}p{11.2cm}@{}}
\toprule
\textbf{Extensiones} &
\begin{itemize}[leftmargin=*, itemsep=2pt]
    \item[*a.] (En cualquier momento) Se pierde la conexión con el
    servicio de reconocimiento: el área de video (4) sigue mostrando la
    imagen de la cámara con normalidad y el sistema reconecta
    automáticamente en unos segundos, retomando el reconocimiento sin
    mostrar ningún error al paciente.
    \item[5a.] Ninguna seña alcanza el 70\,\% de confianza o la seña de
    reposo domina el top-3: las barras del panel (6) se actualizan igual,
    pero no se agrega ninguna palabra a la ventana (5).
    \item[10a.] El sistema no logra traducir la secuencia en una oración
    clara: no se agrega ninguna burbuja al panel (7) y la ventana (5)
    simplemente desaparece.
    \item[10b.] El modelo de lenguaje local no está disponible: la
    burbuja del paso 10 muestra directamente la secuencia de palabras
    reconocidas, sin reformularla en una oración natural, y el flujo
    continúa sin interrupciones.
    \item[11a.] El motor de síntesis de voz no está cargado: el paso 11
    se omite y la burbuja del paso 10 queda igual, sin audio asociado.
\end{itemize} \\
\textbf{Requisitos especiales} & RF-01.1 a RF-01.8; RNF-T-01, RNF-T-05, RNF-01.1. \\
\textbf{Frecuencia} &
Continua durante toda la sesión; se repite cada vez que el paciente forma
una idea completa. \\
\bottomrule
\caption{Cursos alternativos --- CU-01 (real)}
\end{longtable}

% ============================================================================
\subsection{CU-04 (real): Responder al Paciente}

Especializa el caso de uso esencial \textbf{CU-04} (ver Entregable~03,
sección 2.6) sobre la Pantalla \textbf{P-03 --- Estación de Recepción}
(Figura~\ref{fig:cur-p03-cu04}). Incluye, como extensión, el caso de uso
real \textbf{CU-05} (Dictar Respuesta por Voz).

\diagrama{04_mockup_p03_recepcion}{\textwidth}{Pantalla P-03, referenciada por CU-04 (real) y CU-05 (real)}{fig:cur-p03-cu04}

\begin{longtable}{@{}p{4.1cm} >{\raggedright\arraybackslash}p{11.2cm}@{}}
\toprule
\textbf{Nombre} & Responder al Paciente (real) \\
\textbf{ID} & CU-04 (real; esencial en Entregable 03) \\
\textbf{Actor principal} & Personal de Recepción o Enfermería \\
\textbf{Personal involucrado e intereses} &
Personal de Recepción o Enfermería: quiere responder rápido, escribiendo
o dictando, sin pasos adicionales. \newline
Paciente: necesita ver la respuesta en su propia pantalla (P-02) sin
demora perceptible. \\
\textbf{Precondiciones} &
El personal abrió la Pantalla P-03 (\texttt{/recepcion}); el panel
(3) muestra el historial compartido (posiblemente vacío). \\
\textbf{Garantías de éxito} &
La respuesta aparece como una nueva burbuja en el panel (3) de P-03 y en
el panel (7) de P-02, casi de inmediato. \\
\bottomrule
\caption{CU-04 (real) — Responder al Paciente}
\end{longtable}

\setlength{\extrarowheight}{5pt}
\begin{longtable}{@{}>{\raggedright\arraybackslash}p{0.47\textwidth} >{\raggedright\arraybackslash}p{0.47\textwidth}@{}}
\toprule
\textbf{Acción del actor} & \textbf{Respuesta del sistema} \\
\midrule
\endfirsthead
\toprule
\textbf{Acción del actor} & \textbf{Respuesta del sistema} \\
\midrule
\endhead
1. El personal lee la última burbuja del panel (3) (mensaje del
paciente). & \\
2. El personal escribe una respuesta en el área de texto (4). &
3. El sistema habilita el botón ``Enviar'' (5) en cuanto el área de
texto (4) deja de estar vacía. \\
4. El personal hace clic en el botón ``Enviar'' (5). &
5. El sistema agrega la respuesta como una nueva burbuja al final del
panel (3), alineada a la derecha, y vacía el área de texto (4). \\
&
6. La misma burbuja aparece, al mismo tiempo, en el panel (7) de la
Pantalla P-02 del paciente. \\
\bottomrule
\caption{Curso normal de los eventos --- CU-04 (real)}
\end{longtable}
\setlength{\extrarowheight}{0pt}

\begin{longtable}{@{}p{4.1cm} >{\raggedright\arraybackslash}p{11.2cm}@{}}
\toprule
\textbf{Extensiones} &
\begin{itemize}[leftmargin=*, itemsep=2pt]
    \item[2a.] El personal prefiere dictar en lugar de escribir:
    continúa en el punto 1 de CU-05 (real), a continuación, y vuelve al
    paso 4 de este caso de uso con el texto ya cargado en (4).
\end{itemize} \\
\textbf{Requisitos especiales} & RF-01.7, RF-02.1; RNF-T-05. \\
\textbf{Frecuencia} &
Cada vez que el personal necesita responder durante la atención. \\
\bottomrule
\caption{Cursos alternativos --- CU-04 (real)}
\end{longtable}

\subsubsection*{CU-05 (real): Dictar Respuesta por Voz --- extensión de CU-04}

Especializa el caso de uso esencial \textbf{CU-05} (ver Entregable~03,
sección 2.7), en el punto de extensión ``ingreso de texto'' de CU-04, sobre
la misma Pantalla P-03 (Figura~\ref{fig:cur-p03-cu04}).

\begin{longtable}{@{}p{4.1cm} >{\raggedright\arraybackslash}p{11.2cm}@{}}
\toprule
\textbf{Nombre} & Dictar Respuesta por Voz (real) \\
\textbf{ID} & CU-05 (real; esencial en Entregable 03) \\
\textbf{Actor principal} & Personal de Recepción o Enfermería \\
\textbf{Precondiciones} &
El navegador soporta \texttt{SpeechRecognition} (Web Speech API); el
botón de micrófono (6) está habilitado. \\
\textbf{Garantías de éxito} &
El texto dictado queda cargado en el área de texto (4), listo para
enviarse por el paso 4 de CU-04 (real). \\
\bottomrule
\caption{CU-05 (real) — Dictar Respuesta por Voz}
\end{longtable}

\setlength{\extrarowheight}{5pt}
\begin{longtable}{@{}>{\raggedright\arraybackslash}p{0.47\textwidth} >{\raggedright\arraybackslash}p{0.47\textwidth}@{}}
\toprule
\textbf{Acción del actor} & \textbf{Respuesta del sistema} \\
\midrule
\endfirsthead
\toprule
\textbf{Acción del actor} & \textbf{Respuesta del sistema} \\
\midrule
\endhead
1. El personal hace clic en el botón de micrófono (6). &
2. El botón (6) cambia a su estado ``escuchando'' y el indicador (7)
muestra ``Escuchando... click para detener''. \\
3. El personal habla. &
4. El sistema va transcribiendo el habla y sumando el texto reconocido
al área de texto (4), a continuación de cualquier texto ya escrito. \\
5. El personal hace clic nuevamente en el botón (6) para detener el
dictado. &
6. El indicador (7) vuelve a ``Click para hablar''. \\
7. El flujo continúa en el paso 4 de CU-04 (real): el personal revisa el
texto en (4) y hace clic en ``Enviar'' (5). & \\
\bottomrule
\caption{Curso normal de los eventos --- CU-05 (real)}
\end{longtable}
\setlength{\extrarowheight}{0pt}

\begin{longtable}{@{}p{4.1cm} >{\raggedright\arraybackslash}p{11.2cm}@{}}
\toprule
\textbf{Extensiones} &
\begin{itemize}[leftmargin=*, itemsep=2pt]
    \item[1a.] El navegador no soporta reconocimiento de voz: el botón
    (6) aparece deshabilitado y el indicador (7) muestra el aviso ``El
    navegador no soporta reconocimiento de voz''; el personal continúa
    escribiendo directamente en (4).
    \item[4a.] Ocurre un error de micrófono: el indicador (7) muestra
    ``Error de micrófono: ...'' y el sistema deja de escuchar.
\end{itemize} \\
\textbf{Requisitos especiales} & RF-02.2; restricción de navegador (Google
Chrome). \\
\bottomrule
\caption{Cursos alternativos --- CU-05 (real)}
\end{longtable}

% ============================================================================
\subsection{CU-06 (real): Capturar Muestra de Seña para Entrenamiento}

Especializa el caso de uso esencial \textbf{CU-06} (ver Entregable~03,
sección 2.9) sobre la Pantalla \textbf{P-04 --- Modo Captura}
(Figura~\ref{fig:cur-p04-cu06}).

\diagrama{04_mockup_p04_captura}{0.9\textwidth}{Pantalla P-04, referenciada por CU-06 (real)}{fig:cur-p04-cu06}

\begin{longtable}{@{}p{4.1cm} >{\raggedright\arraybackslash}p{11.2cm}@{}}
\toprule
\textbf{Nombre} & Capturar Muestra de Seña para Entrenamiento (real) \\
\textbf{ID} & CU-06 (real; esencial en Entregable 03) \\
\textbf{Actor principal} & Desarrollador de IA / Equipo de Proyecto \\
\textbf{Precondiciones} &
El desarrollador hizo clic en ``Ajustes (Modo Captura)'' (callout (9) de
la Pantalla P-02), abriendo la Pantalla P-04. \\
\textbf{Garantías de éxito} &
Una o más muestras quedaron grabadas; el contador entre paréntesis junto
a la etiqueta, en el selector (2), se actualizó. \\
\bottomrule
\caption{CU-06 (real) — Capturar Muestra de Seña para Entrenamiento}
\end{longtable}

\setlength{\extrarowheight}{5pt}
\begin{longtable}{@{}>{\raggedright\arraybackslash}p{0.47\textwidth} >{\raggedright\arraybackslash}p{0.47\textwidth}@{}}
\toprule
\textbf{Acción del actor} & \textbf{Respuesta del sistema} \\
\midrule
\endfirsthead
\toprule
\textbf{Acción del actor} & \textbf{Respuesta del sistema} \\
\midrule
\endhead
1. El desarrollador elige una etiqueta existente haciendo clic en uno de
los chips del selector (2) (o la busca primero con el filtro (4), si hay
más de seis etiquetas). & \\
2. El desarrollador ajusta el control ``Repeticiones'' (9) según cuántas
muestras seguidas quiere grabar. & \\
3. El desarrollador hace clic en ``Empezar'' (10). & \\
&
4. El sistema muestra una cuenta regresiva (6) de 3 a 1 sobre el área de
video (5). \\
&
5. El sistema activa el indicador ``GRABANDO'' (7) y captura 45 frames
de landmarks y blendshapes faciales desde el área de video (5). \\
&
6. El sistema valida y guarda la muestra; si se
configuraron varias repeticiones, el indicador ``MUESTRA \emph{n} /
\emph{N}'' (8) avanza y el sistema repite los pasos 4 a 6. \\
&
7. Al completar todas las repeticiones, el sistema confirma el guardado
y actualiza el contador de muestras junto a la etiqueta activa, visible
bajo el video (``Grabando para: ... (N muestras)'') y en el chip
correspondiente del selector (2). \\
\bottomrule
\caption{Curso normal de los eventos --- CU-06 (real)}
\end{longtable}
\setlength{\extrarowheight}{0pt}

\begin{longtable}{@{}p{4.1cm} >{\raggedright\arraybackslash}p{11.2cm}@{}}
\toprule
\textbf{Extensiones} &
\begin{itemize}[leftmargin=*, itemsep=2pt]
    \item[1a.] La seña todavía no existe: el desarrollador hace clic en
    ``+ Nueva seña'' (3), escribe el nombre (validado contra
    \texttt{\^{}[a-z0-9\_]\{1,32\}\$}) y lo confirma; el sistema la
    agrega al selector (2) y la deja seleccionada como etiqueta activa.
    La carpeta de datos recién se crea al guardar la primera muestra.
    \item[5a.] El desarrollador hace clic en ``Detener'' (el botón (10)
    cambia de rótulo mientras hay una grabación o cuenta regresiva en
    curso): el sistema cancela la repetición actual, no la guarda, y
    vuelve al estado inicial de la Pantalla P-04.
\end{itemize} \\
\textbf{Requisitos especiales} & RF-T-01 a RF-T-04. \\
\textbf{Frecuencia} &
Bajo demanda, cuando se necesita ampliar o reforzar el vocabulario. \\
\bottomrule
\caption{Cursos alternativos --- CU-06 (real)}
\end{longtable}

% ============================================================================
\subsection{CU-09 (real): Otorgar Consentimiento Explícito de Cámara}

Especializa el caso de uso esencial \textbf{CU-09} (ver Entregable~03,
sección 2.10) sobre la Pantalla \textbf{P-05 --- Diálogo de
Consentimiento} (Figura~\ref{fig:cur-p05-cu09}).

\diagrama{04_mockup_p05_consentimiento}{0.75\textwidth}{Pantalla P-05, referenciada por CU-09 (real)}{fig:cur-p05-cu09}

\begin{longtable}{@{}p{4.1cm} >{\raggedright\arraybackslash}p{11.2cm}@{}}
\toprule
\textbf{Nombre} & Otorgar Consentimiento Explícito de Cámara (real) \\
\textbf{ID} & CU-09 (real; esencial en Entregable 03) \\
\textbf{Actor principal} & Paciente \\
\textbf{Precondiciones} &
El paciente abrió la Estación del Paciente; el sistema muestra el
Diálogo P-05 antes de solicitar el permiso de cámara al navegador. \\
\textbf{Garantías de éxito} &
La cámara se activa solo si el botón (5) fue presionado; en caso
contrario, permanece inactiva. \\
\bottomrule
\caption{CU-09 (real) — Otorgar Consentimiento Explícito de Cámara}
\end{longtable}

\setlength{\extrarowheight}{5pt}
\begin{longtable}{@{}>{\raggedright\arraybackslash}p{0.47\textwidth} >{\raggedright\arraybackslash}p{0.47\textwidth}@{}}
\toprule
\textbf{Acción del actor} & \textbf{Respuesta del sistema} \\
\midrule
\endfirsthead
\toprule
\textbf{Acción del actor} & \textbf{Respuesta del sistema} \\
\midrule
\endhead
1. El paciente lee el texto explicativo (3) sobre el uso de la cámara. & \\
2. El paciente hace clic en ``Acepto y continúo'' (5). &
3. El sistema registra el consentimiento como evento de auditoría
anónimo y recién entonces habilita el hook de captura, que solicita el
permiso de cámara al navegador. \\
&
4. El sistema cierra el Diálogo P-05 y muestra la Pantalla P-02, con el
indicador (3) en ``EN VIVO''. \\
\bottomrule
\caption{Curso normal de los eventos --- CU-09 (real)}
\end{longtable}
\setlength{\extrarowheight}{0pt}

\begin{longtable}{@{}p{4.1cm} >{\raggedright\arraybackslash}p{11.2cm}@{}}
\toprule
\textbf{Extensiones} &
\begin{itemize}[leftmargin=*, itemsep=2pt]
    \item[2a.] El paciente hace clic en ``No acepto'' (4): el sistema
    registra el rechazo como evento de auditoría anónimo, la cámara no
    se activa y el diálogo pasa al estado ``declined'' (6)--(7), con el
    botón ``Reconsiderar'' (7) para volver al estado inicial.
    \item[3a.] El navegador deniega el permiso de cámara: el sistema
    informa el error y la Pantalla P-02 no llega a mostrar video.
\end{itemize} \\
\textbf{Requisitos especiales} & RF-03.1. \\
\textbf{Frecuencia} & Una vez por sesión de la Estación del Paciente. \\
\bottomrule
\caption{Cursos alternativos --- CU-09 (real)}
\end{longtable}

% ============================================================================
\subsection{CU-10 (real): Reportar Traducción Incorrecta o Ambigua}

Especializa el caso de uso esencial \textbf{CU-10} (ver Entregable~03,
sección 2.11) sobre las Pantallas P-02 y P-03, mediante el callout de
reporte agregado a cada burbuja de conversación (callout (10) en P-02,
Figura~\ref{fig:cur-p02-cu01}; callout (8) en P-03,
Figura~\ref{fig:cur-p03-cu04}).

\begin{longtable}{@{}p{4.1cm} >{\raggedright\arraybackslash}p{11.2cm}@{}}
\toprule
\textbf{Nombre} & Reportar Traducción Incorrecta o Ambigua (real) \\
\textbf{ID} & CU-10 (real; esencial en Entregable 03) \\
\textbf{Actor principal} & Paciente, Personal de Recepción o Enfermería \\
\textbf{Precondiciones} &
Existe al menos una burbuja del paciente en el panel de conversación de
P-02 (7) o de P-03 (3); el botón de reporte solo aparece sobre mensajes
del paciente. \\
\textbf{Garantías de éxito} &
El reporte queda registrado en la base de datos relacional local
(\texttt{data/lsa\_salud.db}), asociado al texto y a las palabras del
mensaje señalado. El mensaje se marca como tachado en ambas caras de la
estación (P-02 y P-03) con un aviso operativo, y se solicita al paciente
repetir la seña. \\
\bottomrule
\caption{CU-10 (real) — Reportar Traducción Incorrecta o Ambigua}
\end{longtable}

\setlength{\extrarowheight}{5pt}
\begin{longtable}{@{}>{\raggedright\arraybackslash}p{0.47\textwidth} >{\raggedright\arraybackslash}p{0.47\textwidth}@{}}
\toprule
\textbf{Acción del actor} & \textbf{Respuesta del sistema} \\
\midrule
\endfirsthead
\toprule
\textbf{Acción del actor} & \textbf{Respuesta del sistema} \\
\midrule
\endhead
1. El actor hace clic en el botón ``Reportar traducción'' de una burbuja
del paciente (callout (10) en P-02 o (8) en P-03). &
2. El botón pasa a estado ``Enviando...'' y el sistema registra el
reporte con el texto del mensaje, la estación y las palabras reconocidas. \\
&
3. En ambas pantallas (P-02 y P-03), la burbuja del mensaje se actualiza
inmediatamente:
\begin{itemize}[leftmargin=*, itemsep=2pt]
    \item El texto del mensaje queda tachado (\textit{strikethrough}) y con
    borde/fondo ámbar de advertencia.
    \item El botón pasa a la insignia permanente ``Traducción reportada''.
    \item En la pantalla del paciente (P-02) se muestra: \textit{``Traducción
    reportada. Repetí la seña para enviar el mensaje de nuevo.''}.
    \item En la pantalla de recepción (P-03) se muestra: \textit{``Traducción
    reportada como incorrecta. No la tomes como válida: el paciente va a
    repetir la seña.''}.
\end{itemize} \\
\bottomrule
\caption{Curso normal de los eventos --- CU-10 (real)}
\end{longtable}
\setlength{\extrarowheight}{0pt}

\begin{longtable}{@{}p{4.1cm} >{\raggedright\arraybackslash}p{11.2cm}@{}}
\toprule
\textbf{Extensiones} &
\begin{itemize}[leftmargin=*, itemsep=2pt]
    \item[2a.] La solicitud falla por problema de red o error de validación:
    el botón vuelve a estar disponible con el mensaje \textit{``No se pudo
    enviar el reporte. Probá de nuevo.''} al pie de la burbuja; el mensaje
    no se tacha y el actor puede reintentar.
\end{itemize} \\
\textbf{Requisitos especiales} & RF-06.1, RNF-03.1. \\
\textbf{Frecuencia} &
Bajo demanda, cada vez que se detecta una traducción incorrecta o
ambigua. \\
\bottomrule
\caption{Cursos alternativos --- CU-10 (real)}
\end{longtable}

% ============================================================================
\subsection{CU-12 (real): Finalizar Atención y Preparar la Estación para el
Próximo Paciente}

Especializa el caso de uso esencial \textbf{CU-12} (ver Entregable~03) sobre
la Pantalla \textbf{P-03 --- Estación de Recepción}
(Figura~\ref{fig:cur-p03-cu04}, callout (9)). Su realización técnica se
documenta como \textbf{SO-09} / \textbf{DC-09} en el Entregable~08
(Diagramas de Colaboración).

\begin{longtable}{@{}p{4.1cm} >{\raggedright\arraybackslash}p{11.2cm}@{}}
\toprule
\textbf{Nombre} & Finalizar Atención y Preparar la Estación para el Próximo
Paciente (real) \\
\textbf{ID} & CU-12 (real; esencial en Entregable 03) \\
\textbf{Actor principal} & Personal de Recepción o Enfermería \\
\textbf{Personal involucrado e intereses} &
Personal de Recepción o Enfermería: necesita dejar la estación lista para
el próximo paciente sin que quede a la vista ninguna conversación previa.
\newline
Paciente: espera que su conversación no quede expuesta a la siguiente
persona atendida. \\
\textbf{Precondiciones} &
La atención del paciente actual concluyó; el panel (3) de la Pantalla
P-03 muestra la conversación de esa atención. \\
\textbf{Garantías de éxito} &
El panel de conversación queda vacío en las dos caras de la estación
(P-02 y P-03) y no persiste ningún dato de la atención finalizada (RNF-03). \\
\bottomrule
\caption{CU-12 (real) — Finalizar Atención y Preparar la Estación para el
Próximo Paciente}
\end{longtable}

\setlength{\extrarowheight}{5pt}
\begin{longtable}{@{}>{\raggedright\arraybackslash}p{0.47\textwidth} >{\raggedright\arraybackslash}p{0.47\textwidth}@{}}
\toprule
\textbf{Acción del actor} & \textbf{Respuesta del sistema} \\
\midrule
\endfirsthead
\toprule
\textbf{Acción del actor} & \textbf{Respuesta del sistema} \\
\midrule
\endhead
1. El personal hace clic en el botón ``Finalizar atención'' (callout (9)
de la Pantalla P-03). &
2. El sistema vacía de inmediato el panel de conversación (3) de la
Pantalla P-03. \\
&
3. El sistema notifica a la Pantalla P-02, del lado del paciente, por el
canal de sincronización de la estación; esa pantalla vacía a su vez el
panel de conversación (7). \\
&
4. El sistema registra un evento de auditoría anónimo que marca el fin
de la atención, sin datos identificatorios ni contenido de la
conversación. \\
&
5. La estación queda con ambos paneles vacíos, lista para la siguiente
atención. \\
\bottomrule
\caption{Curso normal de los eventos --- CU-12 (real)}
\end{longtable}
\setlength{\extrarowheight}{0pt}

\begin{longtable}{@{}p{4.1cm} >{\raggedright\arraybackslash}p{11.2cm}@{}}
\toprule
\textbf{Extensiones} &
\begin{itemize}[leftmargin=*, itemsep=2pt]
    \item[3a.] La Pantalla P-02 no está abierta en ese momento: al
    abrirse, el panel de conversación (7) ya inicia vacío, porque ningún
    dato de la atención anterior persiste en el sistema.
\end{itemize} \\
\textbf{Requisitos especiales} & RNF-03.1, RF-04.1. \\
\textbf{Frecuencia} &
Una vez por atención, al finalizar cada paciente. \\
\bottomrule
\caption{Cursos alternativos --- CU-12 (real)}
\end{longtable}
